\documentclass[10pt]{article}

\usepackage[letterpaper, margin=0.85in]{geometry}
\usepackage{iftex}
\usepackage{graphicx}
\usepackage{amsmath, amssymb}
\usepackage{url}
\usepackage{xcolor}
\usepackage[font=footnotesize,labelfont=bf,skip=3pt]{caption}
\usepackage{titlesec}
\usepackage{microtype}
\usepackage{float}

\ifXeTeX
  \usepackage{fontspec}
  \usepackage{unicode-math}
  \setmainfont{Times New Roman}
  \setsansfont{Arial}
  \setmonofont{Consolas}
  \setmathfont{Cambria Math}
\fi

\titlespacing*{\section}{0pt}{6pt}{2pt}
\titleformat{\section}{\normalfont\bfseries}{\thesection.}{0.4em}{}
\captionsetup{justification=justified,singlelinecheck=false}

\setlength{\parskip}{2pt}
\setlength{\parindent}{0pt}
\setlength{\textfloatsep}{5pt}
\setlength{\floatsep}{4pt}
\setlength{\abovecaptionskip}{2pt}
\setlength{\belowcaptionskip}{0pt}

\graphicspath{{../results/manuscript_figures/}{../results/figures/}}

\pagestyle{plain}

\begin{document}

\begin{center}
{\Large \textbf{Spacer composition, not motif grammar, drives regulatory signals in DNA sequence models}}
\end{center}

\section*{Background}

The billboard model of enhancer function holds that regulatory activity depends on which transcription factor binding sites are present---their absolute identity and composition---rather than on their relative positions, orientations, or spacing~[1,2].
This is a contrarian position: the prevailing view assumes that motif grammar (positional syntax) carries regulatory meaning, and DNA foundation models are increasingly benchmarked on their sensitivity to such arrangement.
We test the billboard hypothesis directly.
Standard vocabulary-preserving shuffles confound motif order with spacer DNA, making grammar sensitivity indistinguishable from compositional sensitivity.
We resolve that confound with a spacer-factored grammar sensitivity framework (SF-GSI) that isolates arrangement effects from spacer-composition effects and asks: does positional context survive once absolute context is controlled?

\begin{figure}[H]
\centering
\includegraphics[width=0.7\textwidth]{fig1_spacer_confound.pdf}
\caption{\textbf{Absolute context dominates apparent grammar effects.} Spacer composition (absolute context) explains most measured grammar sensitivity; motif rearrangement (positional context) induces substantially smaller prediction changes.}
\label{fig:spacer}
\end{figure}

\section*{Methods}

SF-GSI starts from the grammar sensitivity index, $GSI = \sigma(\hat{y}) / |\mu(\hat{y})|$, computed over vocabulary-preserving motif shuffles, and factors that response into positional context (motif position/orientation) versus absolute context (spacer composition) using matched-spacer perturbations and null-normalized testing.
The design directly operationalizes the billboard hypothesis: if positional context does not matter, then controlling for absolute context should collapse apparent grammar sensitivity to noise.
To our knowledge, this is the first framework to factorially decompose apparent grammar sensitivity into arrangement and composition components in this setting.
We applied SF-GSI to five MPRA datasets (Agarwal, Klein, Inoue, Vaishnav, and Jores), totaling 7,650 model--enhancer measurements: 3 foundation models $\times$ 5 datasets $\times$ 500 enhancers, plus Enformer on 3 human datasets $\times$ 50 enhancers.
The model set spans transformer, state-space, and MPRA-native families, including DNABERT-2, Nucleotide Transformer v2, HyenaDNA, Enformer, and PARM~[4--7]. Code: \url{https://github.com/bryanc5864/GRAMLANG}.
\begin{figure}[H]
\centering
\includegraphics[width=0.7\textwidth]{fig3_gsi_census.pdf}
\caption{\textbf{Statistical correction collapses apparent grammar prevalence.} Calls fall from 100\% under naive thresholding to 8.3\% after null normalization and to 0.17\% (13/7,650) after FDR correction.}
\label{fig:census}
\end{figure}

\section*{Results}

The dominant result is the correction cascade in Fig.~\ref{fig:census}: under naive thresholding, all tested enhancers appear grammar-sensitive, but this falls to 8.3\% after null normalization and to 0.17\% (13/7,650) after FDR correction.
That three-order-of-magnitude collapse is the signature predicted by the billboard model---once absolute context is controlled, the apparent prevalence of positional syntax vanishes.

SF-GSI explains the mechanism.
As shown in Fig.~\ref{fig:spacer}, absolute context (spacer DNA composition) accounts for 78--86\% of measured grammar-sensitivity variance, while compositional features explain 40--80\% of expression-prediction variance.
Positional context---motif-only rearrangements---produces substantially smaller perturbation effects than spacer replacement or GC-shifting perturbations.
These models have learned \textit{what} is present, not \textit{where} it is; the signal attributed to grammar is, in the vast majority of cases, a compositional signal misread as a positional one.

The residual signal after correction is structured rather than random.
Cross-model agreement is concentrated in specific biological contexts, indicating that dataset context matters more than nominal architecture when asking where credible grammar-sensitive candidates are enriched.
Positive controls confirm that the framework does not erase genuine grammar: when spacer sequence is held fixed in synthetic constructs, orientation effects remain strongly detectable across models ($p \le 7.2\times10^{-15}$, with two models reaching numerical underflow).
However, these controlled cases are rare exceptions, not the rule.
Variance decompositions assign only 0--1.6\% of explained variance to arrangement, versus 8--22\% to motif vocabulary---consistent with the billboard prediction that absolute identity, not positional syntax, is the dominant learned feature.
Cross-species transfer is near zero, further suggesting that whatever positional logic models encode is not generalizable.
These findings matter beyond MPRA grammar analysis because benchmark perturbations that fail to control for absolute context can systematically overstate the positional knowledge that regulatory foundation models have acquired.

\section*{Conclusions}

Our results provide direct quantitative support for the billboard model of enhancer regulation as learned by DNA foundation models: these models encode absolute context---what binding sites and sequence features are present---but show negligible encoding of positional context once composition is controlled.
Matched-spacer positive controls confirm that genuine arrangement sensitivity is detectable in principle, but it accounts for at most 1.6\% of explained variance in practice.
Claims of learned regulatory grammar should be made only under perturbation designs that explicitly factor out absolute context; without that control, compositional signals masquerade as positional ones.

\vspace{4pt}
{\scriptsize
\textbf{References:}
[1]~Arnosti \& Kulkarni, \textit{J.~Cell.~Biochem.} 94:890, 2005.\
[2]~Long et~al., \textit{Cell} 167:1170, 2016.\
[3]~Castro-Mondragon et~al., \textit{Nucleic Acids Res.} 52:D142, 2024.\
[4]~Zhou et~al., \textit{bioRxiv} 2023, DNABERT-2.\
[5]~Nguyen et~al., \textit{bioRxiv} 2023, HyenaDNA.\
[6]~Avsec et~al., \textit{Nature Methods} 18:1196, 2021, Enformer.\
[7]~Agarwal et~al., \textit{Science} 380:eadh0559, 2023; Jores et~al., \textit{Nat.~Plants} 7:842, 2021.
}

\end{document}
